
\minitoc

% ========================================================================================================
% Introduction 

L'objectif du cours d'algorithmique est de donner, pour des problèmes "compliqués" des algorithmes de 
résolution dont la durée d'exécution est "raisonnable". 
En effet, pour beaucoup de problème algorithmiques, la solution naïve est rarement faisable à l'échelle 
humaine. Il faut donc construire des algorithmes dont la durée d'exécution est relativement faible. 
L'objectif de ce cours est donc de définir formellement cette notion de complexité et de donner 
des outils mathématiques pour nous permettre de "comparer" la complexité de deux algorithmes. 


% ========================================================================================================
% Mesure de complexité des algorithmes

\section{Mesure de complexité des algorithmes}

Lors de l'exécution d'un algorithme, nous devons au péalable connaître deux de ses caractéristiques : 
\begin{itemize}
    \item le temps qu'il va mettre à s'exécuter, ce que l'on appellera la \emph{complexité temporelle} 
    \item le nombre de cases en mémoire dont il va devoir se servir, ce que l'on appelera la \emph{complexité spatiale}. 
\end{itemize}
Il est évident que ces deux notions dépendront évidement de la taille des données que l'algorithme devra traiter. 

\begin{definition}[Complexité Temporelle]
    Soit un algorithme $A$ et $D_n$ un ensemble de données de taille $n  $ 
\end{definition}


% ========================================================================================================
% Complexité Asymptotique

\section{Complexité Asymptotique}